\subsection{HTTP-Webserver und lokale Browsersteuerung}
Für die Bedienung der WordClock kann eine browserbasierte Schnittstelle genutzt werden.
Die Grundlage dafür bildet das \gls{http}, das in \gls{iot}-Systemen zur Bereitstellung von Diensten über Server eingesetzt werden kann.
\gls{http}-Server dienen als Anwendungskomponenten, die über \gls{tcp}-Verbindungen mit Clients kommunizieren und Anfragen nach dem Request-Response-Prinzip beantworten~\cite{hou2016iotcloud}.

Im Kontext der WordClock übernimmt der Webbrowser die Rolle des Clients, während der Mikrocontroller als Server agiert.
Der Client sendet eine Anfrage an den Server, beispielsweise zum Abruf einer Bedienoberfläche oder von Statusinformationen.
Der Server verarbeitet diese Anfrage und stellt eine entsprechende Antwort bereit.
Diese Antwort enthält sowohl Informationen über den Bearbeitungsstatus als auch gegebenenfalls angeforderte Inhalte.
Durch dieses Kommunikationsprinzip kann eine einfache und strukturierte Interaktion zwischen Benutzer:innen und dem eingebetteten System realisiert werden~\cite{W3Shttp}.

Für die Umsetzung einfacher Bedienfunktionen sind insbesondere GET-Anfragen relevant~\cite{hou2016iotcloud}.
Ein GET-Request dient als Methode zum abruf oder anfordern von Ressourcen von einem \gls{http}-Server~\cite{W3Shttp}~\cite{hou2016iotcloud}.
Im Kontext der WordClock kann dieses Prinzip genutzt werden, um eine \gls{html}-basierte Bedienseite oder Statusinformationen über den Browser abzurufen.
GET-Requests können zusätzlich Name-Wert-Paare innerhalb der \gls{url} übertragen~\cite{W3Shttp}.

Weiterhin sollen sie nicht für sensible Daten verwendet werden und besitzenLängenbeschränkungen~\cite{W3Shttp}.
Für die WordClock ist die Nutzung dennoch geeignet, da nur einfache Steuerwerte wie Ein-Aus-Zustand, Helligkeit und \gls{rgb}-Farbwerte übertragen werden.

Webbasierte Dienste können über \gls{html}, \gls{css} und JavaScript realisiert werden~\cite{hou2016iotcloud}.
Für die WordClock ist dieser Ansatz geeignet, da damit eine einfache grafische Oberfläche zur lokalen Bedienung und Überwachung umgesetzt werden kann.
Die Bedienung erfolgt plattformunabhängig über einen Webbrowser und ist damit nicht an ein bestimmtes Endgerät gebunden.
Die implementierte einer \gls{http}-Schnittstelle ermöglicht somit eine netzwerkbasierte Steuerung des Systems über einen Webbrowser.
Damit werden grundlegende Voraussetzungen für den Einsatz im Smart-Home-Umfeld geschaffen.